\scanpage{110}
\mspage{59}
V Sommes amalgamées et carrés $W$-cocartésiens dans \inserted{$\mathrm{Cat}$}

\uline{1}. Généralités sur $U\overset{i}{\lhook\joinrel\relbar}X\overset{j}{\relbar\joinrel\rhook}T$, et les sommes amalgamées $\unsure{X'\amalg_U U'|X'\amalg_U T'}$.

1.1. \inserted{\uline{Données de recollement} entre un ouvert et \deleted{son complémentaire} un fermé :} \deleted{Soit} $X$ dans $\mathrm{Cat}$, $U$ un ouvert de $X$, $T$ le complémentaire fermé.
[ \uline{NB} les ouverts correspondent bien aux ouverts de $\mathrm{Ob}\,X$, topologisé \unsure{pour|par} \unsure{la|sa} relation de \deleted{\unsure{pré}} \inserted{préordre} naturelle $x\leq y\Longleftrightarrow\mathrm{Hom}_X(x,y)\ne\varnothing$, de $=$ pour les fermés.
Les topologies associées aux relations de préordre sont celles où toute intersection d'ouverts est encore ouverte, i.e. toute réunion de fermés est encore fermée. Donc les ouverts d'une telle topologie sont les fermés d'une topologie duale, correspondant à la relation de préordre opposée. La relation de préordre associée à la topologie est
\[
x\leq y\Longleftrightarrow x\in\overline{\{y\}},
\]
la topologie associée au préordre \inserted{sur $I$} est celle pour laquelle les ouverts sont les réunions d'ens de la forme $I_{/x}$ (lesquels forment une base d'ouverts), \inserted{ou encore les ``\uline{cribles}'' ($x\in U,y\leq x\Rightarrow y\in U$),} dont les \unsure{fermés} sont les réunions d'ens $ {}_{x\backslash}I$, c'est-à-dire les cocribles.
Pour \unsure{tt} partie \unsure{$A$} \illegible{1} de $I$, il y a un plus petit ouvert contenant $A$, qui est aussi le plus petit voisinage, c'est la réunion des \deleted{\unsure{$S$}} $I_{/x}$ pour $x\in A$. Et dualement il y a un plus plus grand fermé contenu dans $A$, \deleted{\unsure{d}}dont le complémentaire est le plus petit ouvert contenant $B=I\backslash A$, donc $\bigcup_{x\unsure{\in} B}I_{/x}$. L'adhérence de $A$ \unsure{est} la réunion des adhérences \deleted{de ses} $ {}_{x\backslash}I$ de ses pts. ]
% archive sheet number: 110
